% ============================================================================
% INTRODUCCIÓN
% ============================================================================
% La introducción es la presentación clara, breve y precisa del contenido
% del trabajo. NO debe incluir resultados ni conclusiones.
%
% Según la norma UTEM, debe comprender:
%   a) Las razones que motivaron la elección del tema
%   b) La justificación: fundamentos y relevancia del trabajo
%   c) Formulación clara del problema, objetivos generales y naturaleza del trabajo
%   d) Un resumen de cómo se lograron los objetivos planteados
%   e) Una orientación al lector de cómo se ha organizado el texto
% ============================================================================

\chapter*{INTRODUCCIÓN}
\addcontentsline{toc}{chapter}{INTRODUCCIÓN}  % Agrega la Introducción a la tabla de contenido

% --- Contexto y motivación ---
La Universidad Tecnológica Metropolitana (UTEM), en su búsqueda constante por la modernización de sus procesos infraestructurales y la optimización de los recursos destinados a la gestión de sus campus, enfrenta desafíos significativos en la administración de sus espacios de estacionamiento. En el Campus Ñuñoa, específicamente en el área destinada a los funcionarios, el proceso de ingreso vehicular actual depende de un sistema de validación manual que ha demostrado ser ineficiente desde una perspectiva logística y operativa.

Actualmente, el acceso se rige por un protocolo en el que el funcionario debe contactar directamente al personal de conserjería para solicitar la apertura del portón, ya sea mediante llamadas telefónicas o señales visuales. Este modelo obliga al personal a dedicar una parte sustancial de su jornada laboral a tareas de baja complejidad técnica, como es el reconocimiento visual y la apertura de portones, lo que distrae su atención de labores críticas de control y asistencia en el resto de las instalaciones.

% --- Problema de investigación y Justificación ---
Por tanto, el problema central se identifica como una ineficiencia logística crítica en el control de acceso vehicular del Campus Ñuñoa, caracterizada por la dependencia del factor humano para tareas repetitivas, la generación de latencias operativas y la ausencia de una infraestructura digital que sustente la toma de decisiones basada en datos. La solución propuesta busca automatizar este proceso mediante el uso de inteligencia artificial, específicamente detección de placas patentes, para transformar el acceso en un procedimiento autónomo, seguro y eficiente.

% --- Objetivos ---
\section*{Objetivo General}
\addcontentsline{toc}{section}{Objetivo General}
Validar la viabilidad técnica y operativa de un sistema de control de acceso vehicular automatizado mediante el desarrollo de un prototipo funcional basado en modelos de inteligencia artificial para el reconocimiento de patentes, orientado a optimizar la eficiencia y trazabilidad en el ingreso de funcionarios al Campus Ñuñoa de la UTEM.

\section*{Objetivos Específicos}
\addcontentsline{toc}{section}{Objetivos Específicos}
\begin{itemize}
    \item Analizar el estado del arte de las tecnologías de visión computacional y modelos de aprendizaje profundo aplicados a la detección de objetos en tiempo real en entornos de seguridad física.
    \item Diseñar una arquitectura de software basada en el paradigma de \textit{Edge Computing} que permita el procesamiento local de datos para garantizar bajas latencias en la apertura del portón y una mayor protección de la privacidad de los datos.
    \item Entrenar y validar un modelo de detección de placas patentes chilenas que reconozca los cinco formatos definitivos, abarcando desde el sistema clásico de 1985 hasta las nuevas configuraciones aprobadas recientemente.
    \item Implementar una plataforma web unificada que actúe como panel de administración para el registro de funcionarios, y que integre el flujo de captura de video por cámara web local para simular la validación de patentes en tiempo real.
\end{itemize}

% --- Alcances y Limitaciones ---
\section*{Alcances y Limitaciones}
\addcontentsline{toc}{section}{Alcances y Limitaciones}
\label{sec:alcances_limitaciones}
El presente trabajo define una hoja de ruta técnica que abarca desde la investigación conceptual hasta la entrega de un artefacto de software funcional en modo de simulación. Para enmarcar de forma precisa el impacto del proyecto, se establecen las siguientes limitaciones:

\begin{itemize}
    \item \textbf{Exclusión de Motocicletas:} Dado que las normativas chilenas establecen que las motocicletas solo portan placa trasera, y considerando que la geometría de captura de las cámaras en porterías de acceso institucional es inherentemente frontal, resulta físicamente inviable su detección. Las motocicletas quedan excluidas del alcance funcional automatizado. No obstante, a nivel de diseño de hardware, el sistema contempla una topología paralela que mantiene operativa la apertura manual (vía control remoto) para resolver este acceso sin fricción.
    \item \textbf{Vehículos Especiales:} El modelo omitirá deliberadamente del entrenamiento aquellas patentes de carácter especial (ej. Bomberos, Carabineros, diplomáticos o de carga pesada), dado que su flujo de acceso no corresponde al perfil de usuario funcionario definido para este recinto. Al igual que con las motocicletas, el acceso contingente para estos vehículos queda plenamente garantizado por el sistema manual preexistente operado por conserjería.
    \item \textbf{Arquitectura de Referencia vs. Prototipo de Validación:} Si bien el diseño arquitectónico de grado industrial propuesto para la institución (Arquitectura de Referencia) recomienda el uso de placas de borde con aceleración gráfica (NVIDIA Jetson) instaladas físicamente en la portería, el alcance material de este Trabajo de Título se enmarca en la creación de una Prueba de Concepto (\textit{Proof of Concept}). Por ende, el sistema entregable consiste en un Prototipo de Validación en software, el cual será ejecutado en un entorno de computadora personal (PC). Este alcance permite demostrar rigurosamente la viabilidad lógica de la arquitectura y la precisión del modelo de IA, dejando la migración a hardware embebido y la instalación física como trabajo futuro institucional.
    \item \textbf{Topología Distribuida (Nube/Local):} La Arquitectura de Referencia contempla una base de datos y panel de gestión en la Nube (para alta disponibilidad administrativa) y un nodo de inferencia operando localmente (\textit{Edge Computing}) en la portería, garantizando tiempos de respuesta $<100$ ms. Para efectos del Prototipo de Validación, el panel administrativo será desplegado en un entorno de nube gratuito (ej. Render), mientras que el nodo de inferencia se ejecutará de forma estrictamente local (\textit{Localhost}) en la PC de simulación, aislando el procesamiento de video de la latencia de internet.
    \item \textbf{Simulación de Actuadores Físicos:} La interacción con el portón será una simulación digital dentro de la aplicación web. El prototipo mostrará indicadores visuales de "Apertura" o "Cierre" en pantalla, sin ejecutar la instalación de hardware físico (motores, cremalleras ni cableado eléctrico) en el Campus Ñuñoa.
\end{itemize}

% --- Estructura del documento ---
\section*{Estructura del Documento}
\addcontentsline{toc}{section}{Estructura del Documento}
El presente trabajo de título se encuentra estructurado en los siguientes capítulos: 
El \textbf{Capítulo 1} aborda el marco teórico y el estado del arte de las tecnologías de visión computacional y redes neuronales convolucionales. 
El \textbf{Capítulo 2} expone el diseño metodológico de la solución y la arquitectura de infraestructura basada en Edge Computing. 
El \textbf{Capítulo 3} detalla el diseño del software de la solución y la especificación de la infraestructura física y hardware del Nodo de Borde.
El \textbf{Capítulo 4} expone el plan de validación y pruebas técnicas diseñado para evaluar las métricas del modelo de IA y el rendimiento del sistema en tiempo real.
El \textbf{Capítulo 5} presenta el estudio de factibilidad económica, analizando los costos de inversión (CAPEX), costos operativos (OPEX), y evaluando la rentabilidad del proyecto mediante los indicadores de VAN y TIR.
Finalmente, se exponen las \textbf{Conclusiones}, donde se evalúa el cumplimiento de los objetivos planteados y se proponen proyecciones para trabajos futuros.
